% GENERATED FILE. Edit canonical lessons, then run npm run generate:books.
% canonical-source-hash: fnv1a64:31dfaf7495b20510
% canonical-lessons: JA-C335-dantai, JA-C335-kaiin, JA-C335-kaihi, JA-C335-jibun, JA-C335-tanin

\chapter{Groups and Members}
\label{ch:ja-groups-and-members}
\hlchaptermodality{335}

\begin{chapteropening}
\textbf{By the end of this chapter:} \emph{I can say why I join a group or leave it.}

Complete the last lesson of chapter 335: I can say why I join a group or leave it.
\end{chapteropening}

\section[\texorpdfstring{\ja{だんたい} (dantai)}{だんたい (dantai)}]{\ja{だんたい} (dantai) — a group; an organization}

\label{lesson:JA-C335-dantai}



\begin{quote}
Before the new one: say the Japanese for looking after, then the Japanese for to handle.
\end{quote}

\subsection*{You'll want to know: \ja{だんたい}}
\textbf{\ja{だんたい}} — \emph{dantai} — \textquotedblleft{}a group; an organization\textquotedblright{}.

\begin{quote}
\textbf{\ja{だんたいで} \ja{いったら、やすく} \ja{なります。}} — \emph{dantai de ittara, yasuku narimasu} — \textquotedblleft{}If you go as a group, it is cheaper.\textquotedblright{}
\end{quote}

\begin{cousinweb}
\textbf{\ja{だんたい}} (\emph{dantai}) is Sino-Japanese (kango): two Chinese characters read with their on'yomi, readings borrowed from Chinese. They are \emph{dan}, \textquotedblleft{}group\textquotedblright{} (Middle Chinese \emph{dwan}) and \emph{tai}, \textquotedblleft{}body\textquotedblright{} (\emph{thejX}) (KANJIDIC2; Middle Chinese in Baxter's transcription).
\end{cousinweb}

\begin{grammarlens}[title={What you've built}]
One more word for giving reasons and opinions.
\end{grammarlens}

\subsection*{Your turn}
\begin{itemize}
\raggedright
  \item \emph{Say it:} \emph{dantai}
  \item \emph{Say it:} \emph{dantai}, once more
  \item \emph{Say it:} \emph{dantai de ittara, yasuku narimasu} — If you go as a group, it is cheaper.
  \item \emph{Recall:} say the Japanese for looking after, then the Japanese for to handle, then say \emph{dantai} again
\end{itemize}

\begin{tcolorbox}[breakable,colback=teal!4,colframe=teal!35!black,title={Before you move on}]
What does \ja{だんたい} mean? (\textquotedblleft{}A group; an organization\textquotedblright{}.) Say it once more.
\end{tcolorbox}
\section[\texorpdfstring{\ja{かいいん} (kaiin)}{かいいん (kaiin)}]{\ja{かいいん} (kaiin) — a club member; a member}

\label{lesson:JA-C335-kaiin}



\begin{quote}
Before the new one: say the Japanese for to handle, then the Japanese for an organization.
\end{quote}

\subsection*{You'll want to know: \ja{かいいん}}
\textbf{\ja{かいいん}} — \emph{kaiin} — \textquotedblleft{}a club member; a member\textquotedblright{}.

\begin{quote}
\textbf{\ja{かいいんに} \ja{なったら、りょうきんが} \ja{やすく} \ja{なります。}} — \emph{kaiin ni nattara, ryōkin ga yasuku narimasu} — \textquotedblleft{}If you become a member, the fee goes down.\textquotedblright{}
\end{quote}

\begin{cousinweb}
\textbf{\ja{かいいん}} (\emph{kaiin}) is Sino-Japanese (kango): two Chinese characters read with their on'yomi, readings borrowed from Chinese. They are \emph{kai}, \textquotedblleft{}meeting\textquotedblright{} (Middle Chinese \emph{hwajH}) and \emph{in}, \textquotedblleft{}member\textquotedblright{} (\emph{hjwen}) (KANJIDIC2; Middle Chinese in Baxter's transcription).
\end{cousinweb}

\begin{grammarlens}[title={What you've built}]
One more word for giving reasons and opinions.
\end{grammarlens}

\subsection*{Your turn}
\begin{itemize}
\raggedright
  \item \emph{Say it:} \emph{kaiin}
  \item \emph{Say it:} \emph{kaiin}, once more
  \item \emph{Say it:} \emph{kaiin ni nattara, ryōkin ga yasuku narimasu} — If you become a member, the fee goes down.
  \item \emph{Recall:} say the Japanese for to handle, then the Japanese for an organization, then say \emph{kaiin} again
\end{itemize}

\begin{tcolorbox}[breakable,colback=teal!4,colframe=teal!35!black,title={Before you move on}]
What does \ja{かいいん} mean? (\textquotedblleft{}A club member; a member\textquotedblright{}.) Say it once more.
\end{tcolorbox}
\section[\texorpdfstring{\ja{かいひ} (kaihi)}{かいひ (kaihi)}]{\ja{かいひ} (kaihi) — a membership fee}

\label{lesson:JA-C335-kaihi}



\begin{quote}
Before the new one: say the Japanese for an organization, then the Japanese for a club member.
\end{quote}

\subsection*{You'll want to know: \ja{かいひ}}
\textbf{\ja{かいひ}} — \emph{kaihi} — \textquotedblleft{}a membership fee\textquotedblright{}.

\begin{quote}
\textbf{\ja{かいひが} \ja{たかいので、やめました。}} — \emph{kaihi ga takai node, yamemashita} — \textquotedblleft{}The membership fee is high, so I quit.\textquotedblright{}
\end{quote}

\begin{cousinweb}
\textbf{\ja{かいひ}} (\emph{kaihi}) is Sino-Japanese (kango): two Chinese characters read with their on'yomi, readings borrowed from Chinese. They are \emph{kai}, \textquotedblleft{}meeting\textquotedblright{} (Middle Chinese \emph{hwajH}) and \emph{hi}, \textquotedblleft{}expense\textquotedblright{} (\emph{phj+jH}) (KANJIDIC2; Middle Chinese in Baxter's transcription).
\end{cousinweb}

\begin{grammarlens}[title={What you've built}]
One more word for giving reasons and opinions.
\end{grammarlens}

\subsection*{Your turn}
\begin{itemize}
\raggedright
  \item \emph{Say it:} \emph{kaihi}
  \item \emph{Say it:} \emph{kaihi}, once more
  \item \emph{Say it:} \emph{kaihi ga takai node, yamemashita} — The membership fee is high, so I quit.
  \item \emph{Recall:} say the Japanese for an organization, then the Japanese for a club member, then say \emph{kaihi} again
\end{itemize}

\begin{tcolorbox}[breakable,colback=teal!4,colframe=teal!35!black,title={Before you move on}]
What does \ja{かいひ} mean? (\textquotedblleft{}A membership fee\textquotedblright{}.) Say it once more.
\end{tcolorbox}
\section[\texorpdfstring{\ja{じぶん} (jibun)}{じぶん (jibun)}]{\ja{じぶん} (jibun) — oneself}

\label{lesson:JA-C335-jibun}



\begin{quote}
Before the new one: say the Japanese for a club member, then the Japanese for a membership fee.
\end{quote}

\subsection*{You'll want to know: \ja{じぶん}}
\textbf{\ja{じぶん}} — \emph{jibun} — \textquotedblleft{}oneself\textquotedblright{}.

\begin{quote}
\textbf{\ja{じぶんで} \ja{きめた} \ja{ほうが} \ja{いいです。}} — \emph{jibun de kimeta hō ga ii desu} — \textquotedblleft{}You had better decide for yourself.\textquotedblright{}
\end{quote}

\begin{cousinweb}
\textbf{\ja{じぶん}} (\emph{jibun}) is Sino-Japanese (kango): two Chinese characters read with their on'yomi, readings borrowed from Chinese. They are \emph{ji}, \textquotedblleft{}oneself\textquotedblright{} (Middle Chinese \emph{dzijH}) and \emph{bun}, \textquotedblleft{}part\textquotedblright{} (\emph{pjun}) (KANJIDIC2; Middle Chinese in Baxter's transcription).
\end{cousinweb}

\begin{grammarlens}[title={What you've built}]
One more word for giving reasons and opinions.
\end{grammarlens}

\subsection*{Your turn}
\begin{itemize}
\raggedright
  \item \emph{Say it:} \emph{jibun}
  \item \emph{Say it:} \emph{jibun}, once more
  \item \emph{Say it:} \emph{jibun de kimeta hō ga ii desu} — You had better decide for yourself.
  \item \emph{Recall:} say the Japanese for a club member, then the Japanese for a membership fee, then say \emph{jibun} again
\end{itemize}

\begin{tcolorbox}[breakable,colback=teal!4,colframe=teal!35!black,title={Before you move on}]
What does \ja{じぶん} mean? (\textquotedblleft{}Oneself\textquotedblright{}.) Say it once more.
\end{tcolorbox}
\section[\texorpdfstring{\ja{たにん} (tanin)}{たにん (tanin)}]{\ja{たにん} (tanin) — other people; a stranger}

\label{lesson:JA-C335-tanin}



\begin{quote}
Before the new one: say the Japanese for a membership fee, then the Japanese for oneself.
\end{quote}

\subsection*{You'll want to know: \ja{たにん}}
\textbf{\ja{たにん}} — \emph{tanin} — \textquotedblleft{}other people; a stranger\textquotedblright{}.

\begin{quote}
\textbf{\ja{たにんの} \ja{かさを} \ja{つかわないで} \ja{ください。}} — \emph{tanin no kasa o tsukawanaide kudasai} — \textquotedblleft{}Please do not use other people's umbrellas.\textquotedblright{}
\end{quote}

\begin{cousinweb}
\textbf{\ja{たにん}} (\emph{tanin}) is Sino-Japanese (kango): two Chinese characters read with their on'yomi, readings borrowed from Chinese. They are \emph{ta}, \textquotedblleft{}other\textquotedblright{} (Middle Chinese \emph{tha}) and \emph{nin}, \textquotedblleft{}person\textquotedblright{} (\emph{nyin}) (KANJIDIC2; Middle Chinese in Baxter's transcription).
\end{cousinweb}

\begin{grammarlens}[title={What you've built}]
One more word for giving reasons and opinions.
\end{grammarlens}

\subsection*{Your turn}
\begin{itemize}
\raggedright
  \item \emph{Say it:} \emph{tanin}
  \item \emph{Say it:} \emph{tanin}, once more
  \item \emph{Say it:} \emph{tanin no kasa o tsukawanaide kudasai} — Please do not use other people's umbrellas.
  \item \emph{Recall:} say the Japanese for a membership fee, then the Japanese for oneself, then say \emph{tanin} again
\end{itemize}

\begin{tcolorbox}[breakable,colback=teal!4,colframe=teal!35!black,title={Before you move on}]
What does \ja{たにん} mean? (\textquotedblleft{}Other people; a stranger\textquotedblright{}.) Say it once more.
\end{tcolorbox}
